\chapter*{Brev til guten som ville bli til lys}
\addcontentsline{toc}{chapter}{Brev til guten som ville bli til lys}
\markboth{Brev til guten}{}

\begin{center}\itshape
Dette er eit brev til eit fag, skrive som eit brev til ein gut,\\ so du som les kan kjenne deg att i båe.
\end{center}
\vspace{6mm}

\begin{flushright}\itshape\small Frå bakerste rad, ein vårdag\end{flushright}

\noindent Til deg som drog til byen for å lære deg å trylle,

Du tok bussen inn, og rundt deg sat forretningsmennene og åt avisene sine. Du lo av dei. Du visste ikkje då at det var dei du skulle tene, og at faget du var på veg for å lære, ein dag skulle levere dei den same hungeren reinskoren: å ete merksemd, side for side, til det ikkje var meir att. Alt du seinare gjorde, vart ført i ei bok med éi spalte. Det som selde, stod der. Det som vart ete opp, stod ingen stad.

På skulen lærte dei deg formlar. Effektivisering for oppvaskmaskiner, like fram å rekne, umogleg å ta feil av. Dei kalla det grunnlag. Men ein formel som berre kan ha rett, lærer deg ingenting; han fortel deg kva bås tingen står i, ikkje kvifor båsane finst. Dei lærte deg å gjere ei form betre langs den eine aksen nokon hadde bede om, og dei lærte deg aldri det du var komen for: korleis du skulle vite kva som stemmer. Mellom rekninga og valet stod det eit tomt steg, og i det steget stod du åleine, slik dei sjølve hadde stått.

So gjorde du oppvaskmaskina levande, og dei kasta deg ut. Tenk på det. Det var det næraste du kom noko sant, og det var nett det dei ikkje hadde ord for. Eit fag som berre kan måle daudt stoff, har inga bok å føre det levande i. Den dagen materialet sjølv reiser seg og får sin eigen vilje, og den dagen kjem, står faget att med eit blikk og utan ein teori om kva ein agent er. Du gav ei maskin ein vilje. Dei såg ein gut som hadde brote ein regel, og ikkje ein gut som hadde peika rett på holet i faget.

På biblioteket møtte du Prinsessa. Ho var komen for å hjelpe folket, men folket ville ikkje ha hjelp, og kollegene ville det minst av alt. Slik har det alltid vore med den gode meininga i dette faget: ho vert lova på nytt i kvar bølgje, frå dei som ville berge handa frå maskina til dei som ville gjere kvardagen vakker for alle, og ho vert aldri teken imot, fordi ho aldri har eit feste. Ein vil hjelpe, men ingen kan seie nei på vegner av faget, og ein godvilje utan eit nei er berre ei kjensle. Difor rann ho deg or hendene til slutt, som vatn.

De gjekk på bar og synte kvarandre magi. Det var fint ei stund. Men sjå kva rommet var: dei dyktige som nikka til dei dyktige og kalla nikket dom. Prisar, blad, eit nikk frå den som alt var nikka til. Ingen spurde kvar kunsten kvilte; det heldt at ho sat. Eg har sete i mange slike rom. Dei liknar meir enn nokon vil høyre på ein annan salong, der velkledde menn ein gong kjende på skallar og las karakter, og var like visse, og hadde like fullt eit apparat, og tok like fullt feil utan å vite det.

Prinsessa sa at om du ville meistre magien, måtte du bli læregut hjå ein oldmeister. Du sa du skulle finne ein. Du gjorde det aldri. Du sat på bar og trylla blomar i drinkane til folk. «Det lærer du i praksis,» seier dei, og det er den vakraste måten eit fag har funne for å seie at det ikkje kan grunngje seg. Du kunne korleis. Du visste aldri at. Og i staden for å gå og leite etter botnen, pynta du overflata, blome etter blome, til applaus.

For du var ein hamskiftar, og du ville bli til reint lys. Det fyrste ein må kunne, sa du sjølv, er å sjå til botnen i ein ting: stire på meitemarken i dagar, til ein ser inn i, og til sist botnen på, han. Du sa det utan å vite at du sa heile saka. Botnen var grunnen. Han fanst. Ein stol er ein posisjon i eit rom; eit skjel teikna seg som tre tal alt i 1966, lenge før faget ditt var ferdig født; meitemarken med. Tre andre fag hadde lagt delane ferdige på hylla, kvar sin, medan faget ditt såg ein annan veg og kalla botnen for fin til å fangast. Du måtte ikkje bli til lys for å sjå. Du måtte berre måle i staden for å døype.

Men du stira ikkje på botnen. Du stira på ei lyspære. Og lyspæra åt deg. Han som ville bli til lyskjegle, til den vide horisonten som rommar kvar form som kan tenkjast, fall for den eine pæra som rommar berre seg sjølv. Demonen åt av det inste ditt, fordi du ikkje hadde noko nei å verje deg med, og då han hadde overtaket, var du ikkje lenger din eigen. Det er ikkje ei segn. Det er eit forretningskart: forma fangar agenten, agenten gjev merksemd og data, og data lukkar sirkelen att til forma. Du bygde sløyfa og gjekk i henne.

Prinsessa flytta heim til slottet. Du prøvde å bli kvitt demonen og greidde det ikkje, for det finst ikkje noko organ i faget ditt som kan støyte ut. Ingen kan ekskluderast frå design, fordi det ikkje finst nokon standard å bli ekskludert frå. Det du helst ville kaste ut, eigde deg; og det einaste du visst kunne miste, den gode meininga, gjekk si veg.

So vart du forvist frå barane og drog ut or byen. Du møtte rare folk og gjekk gjennom syke landskap. Det einaste faget ditt nokon gong greidde å sende frå seg med glans, sende du frå deg her: ein tom prosess kven som helst kunne ta i bruk, nett av di han ikkje kravde fagkunnskap. Eit fag med ein grunn held på avkomet sitt. Du såg ditt drive heimanfrå og slå rot der det fyrst fanst ei ramme å feste det i.

Til sist kom du til slutten, der alt var kvitt. Det kvite var grunnløysa sjølv, og i staden for å stå i henne, trylla du fram duskregn og ein heil verd av regnbogar. Det er den finaste evna du har: å lage ei meining so vakker at ho aldri kan ta feil. Det vrange problemet som per definisjon ikkje har eit rett svar; den tause kunnskapen ingen observasjon kan motseie. Ein regnboge kan ikkje fellast. Difor lærer han deg ingenting. Du stod på ingen grunn og kalla det utsikt.

Du vandra gjennom ruinane av di eiga fortid. Gløymde eignelutar låg strødde i sanden, fjell av ting laga for å kastast. Tårn av utsette planar reiste seg mot skyene, kvar plan eit fundament lova og aldri bygd. Og elva rann, full av augneblink du gjerne ville teke att: folk du køyrde over, folk du svik, folk som snudde deg ryggen. Du stakk handa i vatnet for å trøyste dei, men det var for seint, og det gylne håret til Prinsessa rann deg or fingrane. Det er den andre spalta i boka di, den du aldri førte. No renn ho forbi deg, heile på ein gong.

Lenger borte stod brua, pynta med statuar av deg sjølv. Du ved pulten, ivrig med handa i veret. Du på veg opp ein lyktestolpe. Du som matar Prinsessa med pannekaker, den finaste. So vart dei verre. Du slegen ut på fortauet. Du og Prinsessa som kastar istappar mot kvarandre. Du bøygd over vasken med ein stor kniv i handa. Det er sigerssoga og rekneskapen på same bru: fyrst faget slik det likar å sjå seg sjølv, so det same faget langs den eine aksen det vart dyktig på, ført heilt ut. Kniven over vasken er den same handa som ein gong gjorde maskina levande. Det kjølege vatnet rann nedover kinna på statuane, og det var ikkje du som gret. Det var rekninga som endeleg kom.

Du sprang for å koma deg unna regnet, og sanden vart tung. Det small i himmelen, og du klynka: finst det ingen stad å gøyme seg, har ikkje galningen som bygde dette stelt med tak? Men høyr deg sjølv. Du var galningen. Taket du saknar er grunnen du aldri la, og so lenge du trur at faget er grunnlaust av natur, ber du fråvêret som ein lagnad. Det er det ikkje. I det nokon syner at ein grunn lèt seg gjere, sluttar fråvêret å vere ein lagnad og vert eit val. Du valde det. Kvar dag valde du det om att.

So snubla du ned nokre trapper, i eit hol i bakken, og hamna i ein tunnel. Du nådde ein botn til slutt, men det var ikkje den du leitte etter. På veggen hang ei dame i badedrakt og solhatt, og over midja stod det: «får du nokon gong lyst til å berre koma deg vekk?» Det er reklamen som møter deg på botnen, ikkje grunnen. Det er sløyfa di som lukkar seg: du drog til byen for å sleppe unna, og du lærte eit fag som lever av å selje andre den same flukta. Heilt ned, og framleis berre overflate. Heilt rundt, og attende til forretningsmennene på bussen.

\vignett

Eg skriv ikkje for å la deg liggje i tunnelen. Eg skriv for å seie deg det ingen sa i den salen: det finst ein botn, og han var der heile tida. Form er ein posisjon i eit rom av moglegheiter, forma av samtidige seleksjonstrykk, navigert av agentar på fleire skalaer. Det er heile setninga. Meitemarken du stira på, stolen på kjøkenet, skjelet i tidsskriftet: kvar av dei eitt punkt i eit rom du kunne ha målt, med trykk du kunne ha kjent og ein horisont du kunne ha lese. Det verste er ikkje at botnen var for djup. Det verste er at han var nådd alt, i tre andre fag, og låg og venta medan du trylla blomar.

For elles veit eg kva du er. Du er mannen med fingrane på skallen, viss og dyktig, lenge etter at faget hans er over, fordi ingen enno har opna vevet ved sida av han. Nokon kjem til å opne forma slik Broca opna hjernen, og du kjem ikkje til å kjenne skilnaden den dagen det skjer, og ikkje dagen etter. Med mindre du sjølv reiser deg, slik studenten reiste seg i salen og spurde, roleg, inn i rommet: kva stemmer? Det er det einaste spørsmålet som tel. Neste gong nokon spør det, skal det finnast eit svar i rommet.

So riv. Ikkje av forakt, men av omhug for det som kan stå att. Du ville bli til lys. Bli heller til lyskjegla og ikkje til pæra: til horisonten som rommar formene, ikkje til den eine som åt deg. Grunnen er alt lagd, og han ventar ikkje på løyve. Reis deg or tunnelen. Det er framleis langt att til byen.

\begin{flushright}\itshape frå ein som sat bakerst og noterte\end{flushright}
